\chapter{Conclusions and Future Directions}
\label{Section:Conclusion}

This thesis studies constrained optimization and online algorithms with
costly domain projections, states that depend on past actions, and limited
feedback.

\section{Main results}

Chapter~\ref{Section:projfree-domain} develops a method for nonsmooth convex
optimization that couples linear optimization on the original domain with
projected subgradient updates on an auxiliary set containing it.
With parameters chosen for a query budget $T$, the method attains an
objective gap of $O(T^{-1/2})$ using at most $T$ calls to each of two
oracles: a linear optimization oracle over the original domain and a
subgradient oracle on the auxiliary set. Projections are required only
onto the auxiliary set, which is chosen to have a simple, inexpensive
projection, as with a closed Euclidean ball containing the original domain.
The auxiliary set may also be $\Real^n$, where projection is the identity,
provided the required subgradient access and uniform bound hold globally.
The returned point belongs to the original domain.

Chapter~\ref{Section:projfree} applies the algorithm of Asgari and
Neely~\cite{asgari2024nonsmooth} to the positive parts of the constraints.
Keeping the functional inequalities outside the linear oracle preserves
the tractability of the oracle on the original domain.
With parameters chosen for a query budget $T$, it gives an $O(T^{-1/2})$
upper bound on the objective gap at the returned domain point.
The constraint-violation norm has the same rate if a Lagrange
multiplier exists. The algorithm need not know the multiplier's value,
although its norm enters the violation bound.

Chapter~\ref{Section:ChapEnergyharvest} studies how to allocate energy from a
battery among $n$ subsystems when the battery receives random energy arrivals.
The update rule separates total usage from allocation
proportions: battery feedback regulates expenditure, while entropic mirror
descent updates the proportions on a fixed simplex.
The benchmark is the best nonnegative fixed allocation
whose total usage does not exceed the mean energy arrival.
Under the stated arrival and information assumptions, tuning for a planned
horizon of $T$ slots gives expected cumulative regret
$O(\sqrt{T(1+\log n)})$ with sufficient battery capacity $O(\sqrt T)$.
The controller uses only the energy available in each slot without knowing
the mean arrival.

The energy-allocation benchmark in Chapter~\ref{Section:ChapEnergyharvest}
satisfies a mean budget and may exceed the available energy in individual
slots. Each SDAC benchmark in Chapter~\ref{Section:ChapControl} must satisfy
the energy constraint in every slot, using the storage level generated by
its own past actions.

Chapter~\ref{Section:ChapControl} studies scalar storage control with arrivals
in $[0,1]$ and full cost feedback. The entropic SDAC controller achieves
$O(\sqrt{T\log(H+1)})$ regret against SDAC
policies with memory $H$. With the stated horizon-dependent memory and step
size, it achieves $O(\sqrt{T\log\log T})$ regret against the fixed
infinite-memory SDAC class.

Theorem~\ref{ch4:thm:sdac-lower-bound} establishes a worst-case
expected-regret lower bound of $\Omega(\sqrt T)$ against the fixed
memory-$H$ SDAC class for every causal feasible controller, including
randomized controllers. This matches the upper bound's horizon dependence
for fixed memory and model parameters. The optimal dependence on memory
length and retention remains open.

Chapter~\ref{Section:ChapRateAdaptation} studies rate adaptation using
success or failure acknowledgments to learn mean offered service across
transmission rates. It assumes a known
lower bound $\varepsilon>0$ on the difference between the best achievable mean
offered service and the mean arrival rate. This supplied bound determines
a finite rate grid containing a stabilizing rate. Starting from an empty queue, the UCB1
controller bounds the time-average expected queue length by
$O((1+\log(1/\varepsilon))/\varepsilon^2)$ uniformly over all horizons.

Theorem~\ref{ch5:lower:theorem}, restated without proof from the original
joint paper~\cite{wijewardena2026banditbased}, gives a lower bound of
$\Omega(\varepsilon^{-2})$ for every causal learning policy on some
admissible instance and horizon.
For our UCB1 controller, Theorem~\ref{ch5:upper:thm:ucb-log-lower} gives a
matching $\Omega(\varepsilon^{-2}\log(1/\varepsilon))$ lower bound at a
horizon $T_\varepsilon=\Theta(\varepsilon^{-3}\log(1/\varepsilon))$,
even with deterministic arrivals and capacities. 

\section{Open problems}

\paragraph{Strongly convex objectives.}
An open question is whether the techniques of
Chapter~\ref{Section:projfree-domain} can exploit strong convexity to
achieve a rate of convergence of the objective gap faster than
$O(T^{-1/2})$ for a query budget $T$. The aim is to retain an ordinary
linear optimization oracle over the original domain and projections only
onto an auxiliary set, without requiring the objective to be smooth.

\paragraph{Constraint dependence and violation lower bounds.}
For Chapter~\ref{Section:projfree}, can the dependence on the number $m$
of functional constraints be improved? The constraint-violation bound depends
on $G$, which combines bounds on the sizes of the constraint subgradients,
and on the norm of the Lagrange multiplier.
A related problem is to establish lower bounds for the constraint-violation
norm $V_T$ under the same oracle model and requirement on objective accuracy,
and thereby determine the best possible dependence on $T$ and $m$.

\paragraph{Feasible policy benchmarks for energy allocation.}
Can the regret guarantees against fixed allocations in
Chapter~\ref{Section:ChapEnergyharvest} be extended to feasible
policies that spend up to a fixed amount each slot? Such a policy
fixes a maximum expenditure $u\in[0,1]$ and allocation proportions,
spends $u$ whenever sufficient energy is available, and otherwise spends
all available energy. Under the chapter's arrival timing, its total usage
is $U_t^\pi\coloneqq\min\{u,B_t^\pi+A_t\}$, where $B_t^\pi$ is its own battery
level. Each policy is evaluated with the same arrivals, capacity, and
initial energy as the controller. The benchmark would select the best
fixed threshold and allocation proportions in hindsight.

\paragraph{Memory and retention dependence in storage control.}
For fixed $L_x,L_u>0$ and horizons satisfying
\eqref{ch4:eq:sdac-lower-horizon}, as $\alpha\uparrow1$ the storage
lower bound above is $\Omega((1-\alpha)^{-1}\sqrt T)$.
The upper bound in Theorem~\ref{ch4:thm:regret-bound} is
$O((1-\alpha)^{-3/2}\sqrt{T\log(H+1)})$ in this regime, by
Remark~\ref{ch4:rem:retention-dependence}.
For fixed-memory policies, the unresolved factors are
$\sqrt{\log(H+1)}$ in memory dependence and $(1-\alpha)^{-1/2}$ in
retention dependence. For infinite-memory policies, it is also unknown
whether the regret guarantee must contain the factor
$\sqrt{\log\log T}$.

\paragraph{Uniform queue bounds.}
This thesis leaves open whether a causal learning policy can achieve
$\overline Q_T=O(\varepsilon^{-2})$ uniformly over all horizons under the
assumptions of Chapter~\ref{Section:ChapRateAdaptation}.
Theorem~\ref{ch5:lower:theorem} establishes a lower bound of this order.
Theorem~\ref{ch5:upper:thm:ucb-log-lower} shows that removing the logarithm
requires a different controller, rather than a sharper analysis of the
current UCB1 policy alone.

\paragraph{Feedback along the executed trajectory.}
The current method in Chapter~\ref{Section:ChapControl} uses the full cost
function to obtain a subgradient at the state and action that would result
from using its current policy parameters from the start. These can differ
from the actual state and action because earlier decisions used different
parameters. An open problem is to obtain similar regret bounds using feedback
only at the states and actions actually visited.

\paragraph{Horizon-independent storage control.}
Can the storage regret rates be attained without knowing the horizon in
advance? Online projected subgradient descent with
$\eta_t\coloneqq t^{-1/2}$ achieves $O(\sqrt T)$ regret without knowing $T$ in
advance on a compact convex domain with uniformly bounded subgradients
\cite[Theorem~1]{zinkevich2003online}. For the scalar storage controller, the
memory length in
\eqref{ch4:eq:def-infinite-memory-length}
also depends on the horizon. A restart scheme must account for the current
storage state and the continuing influence of past arrivals.

